\chapter{El topos de espacios}
\label{espacios}

\begin{corollary}[label=cor:bivaluado]
  El clasificador de subobjetos de $\esp$ tiene dos puntos.
\end{corollary}
\begin{proof}
  Sea $\Omega$ el clasificador de subobjetos de $\esp$, entonces $\sub(1)\cong \homs(1,\Omega)$
  y por la observación anterior, $\sub(1)=\{0,1\}$, así que $\homs(1,\Omega)=\{0,1\}$.
\end{proof}

\begin{definition}
  Denotaremos con $\toposM$ a la subcategoría plena de espacios $A$ tales que $c_A$ es un isomorfismo.
\end{definition}

\begin{proposition}
  Dado cualquier espacio $A$, $\Gamma A$ esta en $\toposM$.
\end{proposition}
\begin{proof}
  Sea $A$ un espacio, queremos demostrar que $c_{\Gamma A}$ es un isomorfismo. Por un lado,
  $c_{\Gamma A}\colon\Gamma\Gamma A\to\Gamma A$ es el único subobjeto discreto de $\Gamma A$
  tal que todo punto de $\Gamma A$ se factoriza a través de $c_{\Gamma A}$, pero $\text{id}\colon\Gamma A\to\Gamma A$
  también satisface esta propiedad, entonces por el axioma \ref{ax:gds7}, $c_{\Gamma A}$ es 
  equivalente a $\text{id}$ como subobjeto, por lo que $\Gamma\Gamma A\cong \Gamma A$ y $c_{\Gamma A}$ es un isomorfismo.
\end{proof}

\begin{theorem}[label=thm:funtordiscretos]
  \ref{ax:gds6} y \ref{ax:gds7} inducen una pareja de funtores adjuntos entre $\toposM$ y $\esp$.
\end{theorem}
\begin{proof}
  Para cada espacio $A$ 
  tenemos un monomorfismo $c_A\colon  \Gamma A \to A$ dado
  que $\Gamma A$ es un subobjeto. Veamos que $c_A$ es couniversal hacia
  $A$ desde $\Delta\colon  \toposM \to \esp$, la inclusión de $\toposM$. Sea $B$ un
  objeto de $\toposM$ y $f\colon  \Delta B \to A$,
  deseamos encontrar un morfismo $\overline{f}\colon  \Delta B \to \Delta\Gamma A$ que haga conmutar al siguiente diagrama:
  \begin{equation*}
    \begin{tikzcd}
      \Delta B \ar{r} \ar[swap]{rd}{f} & \Delta\Gamma A \ar[tail]{d}{c_A} \\
        & A.
    \end{tikzcd}
  \end{equation*}
  Notemos que si $B$ fuera el objeto inicial ya acabamos. En otro caso, 
  consideramos al producto fibrado entre $f$ y $c_A$. Además, por \ref{ax:gds6}, podemos
  tomar $b\colon  1\to B$ y analizar qué sucede con el producto fibrado.
  \begin{equation*}
    \begin{tikzcd}
        1\ar[bend left=30]{rrd}\ar[bend right=30]{rdd}[swap]{b}
          \ar[dashed]{rd}{}\\[-2ex]
        &[-2ex] f^{-1} (\Gamma A)\ar{r}{\overline{f}} \ar[tail]{d}[swap]{\cong} & \Gamma A\ar[tail]{d}{c_A}\\
        & B\ar{r}[swap]{f} & A.
    \end{tikzcd}
  \end{equation*}
  Para ver que el exterior conmuta recordemos que según \ref{ax:gds7}, los puntos de $A$ se factorizan a
  través de $\Gamma A$ y $fb$ es un punto de $A$. Entonces los puntos de $B$ se factorizan
  a través de $f^{-1} (\Gamma A)$. Pero $f^{-1} (\Gamma A)$ es discreto ya que es subobjeto
  de $B$, que es discreto, por lo que $\Gamma B \cong f^{-1} (\Gamma A)$ y dado que $B$ está
  en $M$, $\Gamma B \cong B$. Así que $\overline{f}$ es el morfismo que cumple lo que queríamos
  y además es único gracias a que $c_A$ es mono. Por tanto, $\Gamma\colon \esp \to M$ es
  un funtor que manda a los espacios a sus espacios discretos asociados
  y $\Gamma$ es adjunto derecho de $\Delta$ con 
  $c\colon \Delta\Gamma \to \text{Id}$ la counidad.
  
\end{proof}



\begin{proposition}
  Si $A$ es un espacio discreto, entonces $\Gamma A \cong A$.
\end{proposition}
\begin{proof}
  Por \ref{ax:gds7} existe un único subespacio discreto $\Gamma A$ de $A$ tal que todo punto de $A$ está en $\Gamma A$. 
  Dado que $A$ es discreto, $A$ es un subespacio de $A$, y cada punto de $A$ está en $A$, concluimos que $\Gamma A \cong A$.
\end{proof}

\begin{corollary}[label=cor:unidadidentidad]
  $\Gamma\Delta$ es la identidad en $\toposM$ y $\Delta\Gamma$ es idempotente.
\end{corollary}
\begin{proof}
  
\end{proof}

\begin{lemma}
  $\mathcal{C}=(\Delta\Gamma,c,\cong)$ es una comonada en $\esp$.
\end{lemma}
\begin{proof}
  Para todo espacio $A$, tenemos que $\Delta\Gamma A \cong \Delta\Gamma\Delta\Gamma A$ por \ref{cor:unidadidentidad},
  lo cual se puede pensar como un isomorfismo natural $\Id\colon \Delta\Gamma\to \Delta\Gamma\Delta\Gamma$, cuyas componentes son identidades. En esta prueba 
  $c$ es la counidad de la adjunción $\Delta\dashv\Gamma$, es decir, $c\colon \Delta\Gamma\to \Id_{\esp}$ que construimos en \ref{thm:funtordiscretos}.
  Procedemos a demostrar que se cumplen los diagramas de comonada para $\mathcal{C}$. Primero, la coasociatividad,
  \begin{equation*}
    \begin{tikzcd}
      \Delta\Gamma A\ar{r}{\id_{\Delta\Gamma A}}\ar[swap]{d}{\id_{\Delta\Gamma A}} & \Delta\Gamma\Delta\Gamma A\ar{d}{\id_{\Delta\Gamma\Delta\Gamma A}} \\
      \Delta\Gamma\Delta\Gamma A\ar[swap]{r}{\Delta\Gamma \id_{\Delta\Gamma A}} & \Delta\Gamma\Delta\Gamma\Delta\Gamma A.
    \end{tikzcd}
  \end{equation*}
  El diagrama conmuta gracias a que $\Delta\Gamma$ es functorial. Ahora, la counidad, queremos ver
  que el siguiente diagrama conmuta:
  \begin{equation*}
    \begin{tikzcd}
      {} & \Delta\Gamma A\ar{d}{\id}\ar[swap]{ld}{\id}\ar{rd}{\id} & {} \\
      \Delta\Gamma A & \Delta\Gamma\Delta\Gamma A\ar{l}{c_{\Delta\Gamma A}}\ar[swap]{r}{\Delta\Gamma c_A} & \Delta\Gamma A.
    \end{tikzcd}
  \end{equation*}
  Dado que $u_{\Gamma A} = \id_{\Gamma A}$, sucede que $\id_{\Delta\Gamma A} = \Delta u_{\Gamma A}$, y
  ambos triángulos se sigue de las ecuaciones unidad-counidad, por lo que el diagrama conmuta.

\end{proof}

\begin{theorem}[label=thm:toposdiscretos]
  La categoría $M$ de espacios discretos es un topos.
\end{theorem}
\begin{proof}
  Basta demostrar que $M$ es equivalente al topos $\esp_C$ de coálgebras, con
  $\mathcal{C}$ la comónada del lema anterior. Con tal propósito, proponemos al funtor
  $F\colon M\to \esp_{\mathcal{C}}$ que envía a cada espacio discreto $D$ a la coálgebra 
  $(\Delta D, u_D)$, donde $u\colon \Id\to\Delta\Gamma$ es la unidad de la adjunción, y en 
  morfismos $F$ es la identidad. Primero veamos que $(\Delta D, \Delta u_D)$ realmente es una coálgebra,
  para lo cual es necesario que los siguientes diagramas conmuten.
  \begin{equation*}
    \begin{tikzcd}
      \Delta D & \Delta\Gamma\Delta D\ar[swap]{l}{c_{\Delta D}} & {} & \Delta\Gamma\Delta\Gamma\Delta D & \Delta\Gamma\Delta D\ar[swap]{l}{\Delta\Gamma\Delta u_D} \\
       & \Delta D\ar{lu}{\id}\ar[swap]{u}{\Delta u_D} & {} & \Delta\Gamma\Delta D\ar{u}{\id} & \Delta D.\ar{l}{\Delta u_D}\ar[swap]{u}{\Delta u_D}
    \end{tikzcd}
  \end{equation*}
  El primer diagrama conmuta ya que es una ecuación unidad-counidad, mientras que
  el segundo diagrama se da gracias a que las componentes de la unidad son identidades.
  Además, dejar igual a los morfismos resulta en morfismos de coálgebras ya que si $f\colon D\to D'$
  es un morfismo de espacios discretos, entonces el siguiente diagrama conmuta:
  \begin{equation*}
    \begin{tikzcd}
      \Delta D\ar{r}{f}\ar[swap]{d}{\Delta u_D} & \Delta D'\ar{d}{\Delta u_{D'}} \\
      \Delta\Gamma\Delta D\ar[swap]{r}{\Delta\Gamma f} & \Delta\Gamma\Delta D'.
    \end{tikzcd}
  \end{equation*}
  Recordando que $\Delta$ es la inclusión, tenemos que el diagrama es un cuadrado de naturalidad de $u$.
  Por tanto, $F$ es funtorial ya que no hace nada en morfismos. Ahora, proponemos al funtor inverso $G\colon \esp_{\mathcal{C}}\to M$ 
  que envía a cada coálgebra $(A,s)$ a $\Gamma A$ y a cada morfismo de coálgebras $f\colon (A,s)\to (B,t)$ a $\Gamma f$.
  Es fácil ver que $G$ es un funtor ya que $\Gamma$ lo es. Sea $D$ un espacio discreto y 
  $(A,s)$ una coálgebra, veamos qué sucede con las composiciones.
  \begin{align*}
    GF(D) &= G(\Delta D, \Delta u_D) = \Gamma\Delta D = D \\
    FG(A,s) &= F(\Gamma A) = (\Delta\Gamma A, \Delta u_{\Gamma A}) \cong (A,s).
  \end{align*}
  Para convencernos del último paso consideremos el siguiente diagrama:
  \begin{equation*}
    \begin{tikzcd}
      \Delta\Gamma A\ar[swap]{d}{\Delta u_{\Gamma A}}\ar{r}{c_A} & A\ar{d}{s} \\
      \Delta\Gamma\Delta\Gamma A\ar[swap]{r}{\Delta\Gamma c_A} & \Delta\Gamma A.
    \end{tikzcd}
  \end{equation*}
  Se trata de un diagrama conmutativo puesto que la parte inferior es la identidad, 
  una vez más, al ser una ecuación unidad-counidad. Luego, la condición de counidad para que
  $(A,s)$ sea coálgebra nos dice que $c_As = \id_A$, por lo que $c_A$ es mono y tiene inverso por la derecha,
  así que $c_A$ es de hecho un isomorfismo con $s$ su inverso por ambos lados. Entonces la parte superior
  del diagrama también es la identidad, y por lo anterior, $c_A$ es un isomorfismo de coálgebras.
  Por lo tanto $F$ es una equivalencia de categorías y $M$ es un topos.
\end{proof}

\begin{theorem}
  Si $\Omega$ es el clasificador de subobjetos en $\esp$, entonces $\Gamma\Omega$ con 
  $\Gamma v\colon \Gamma1\to\Gamma\Omega$ es el clasificador de subobjetos de M. $\Gamma$ preserva al 
  clasificador de subobjetos.
\end{theorem}
\begin{proof}
  Para cada $A\in M$ y $m\colon S\to A$ un subespacio de $A$, existe un único morfismo $\chi_m\colon A\to \Omega$ tal que el siguiente diagrama conmuta:
  \begin{equation*}
    \begin{tikzcd}
      \Delta S\ar{r}{!}\ar[swap,tail]{d}{\Delta m} & 1\ar{d}{v} \\
      \Delta A\ar[swap]{r}{\chi_m} & \Omega.
    \end{tikzcd}
  \end{equation*}
  Como $\Gamma$ es adjunto derecho, preserva límites, entonces al aplicar $\Gamma$ al diagrama anterior
  obtenemos un producto fibrado. Por \ref{prop:terminaldiscreto} $\Gamma 1 = 1$ y aplicamos \ref{cor:unidadidentidad}.
  \begin{equation*}
    \begin{tikzcd}
      S\ar{r}{!}\ar[swap,tail]{d}{m} & 1\ar{d}{\Gamma v} \\
      A\ar[swap]{r}{\Gamma \chi_m} & \Gamma\Omega.
    \end{tikzcd}
  \end{equation*}
  Si hubiera otro morfismo $f\colon A\to \Gamma\Omega$ tal que el diagrama sea un producto fibrado, entonces
  consideramos al siguiente diagrama:
  \begin{equation*}
    \begin{tikzcd}
      S\ar{r}{!}\ar[swap,tail]{d}{m} & \Delta\Gamma 1\ar{d}{\Delta\Gamma v}\ar{r}{c_1} & 1\ar{d}{v} \\
      A\ar[shift left=0.75ex]{r}{f}\ar[swap,shift right=0.75ex]{r}{\Gamma \chi_m} & \Delta\Gamma\Omega\ar[swap,tail]{r}{c_{\Omega}} & \Omega.
    \end{tikzcd}
  \end{equation*}
  El cuadrado de la derecha conmuta por la naturalidad de la counidad, pero es también un producto fibrado.
  \begin{equation*}
    \begin{tikzcd}
      X\ar[bend left=30]{rrd}{!}\ar[bend right=30]{rdd}[swap]{p}
          \ar[dashed]{rd}{!}\\[-2ex]
      &[-2ex] \Delta\Gamma 1\ar{r}{c_1}\ar[swap]{d}{\Delta\Gamma v} & 1\ar{d}{v} \\
      &\Delta\Gamma\Omega\ar[tail]{r}{c_{\Omega}} & \Omega.
    \end{tikzcd}
  \end{equation*}
  Bajo la hipótesis de que $v!=c_{\Omega}p$, como los diagramas que terminal en $1$ conmutan, 
  $c_{\Omega}p=vc_1!=c_{\Omega}\Delta\Gamma v!$, entonces $p=\Delta\Gamma v!$ y en efecto, es un producto fibrado.
  De esta manera, por el lema del producto fibrado, el exterior es un producto fibrado, lo cual implica que
  $c_{\Omega}f=\chi_m=c_{\Omega}\Gamma \chi_m$, pero $c_{\Omega}$ es mono, entonces $f=\Gamma \chi_m$. 
  Por tanto, $\Gamma\Omega$ con $\Gamma v$ es el clasificador de subobjetos de $M$.
\end{proof}

\begin{theorem}[label=thm:puntitos]
  Si $A$ tiene una cantidad finita de puntos, entonces $\Gamma A=\sqcup 1$ donde la cantidad de copias
  de $1$ es igual a la cantidad de puntos de $A$.
\end{theorem}
\begin{proof}
  Basta demostrarlo para $A$ con dos puntos, $a,b\colon 1\to A$ con $a\neq b$ y el resultado se sigue por inducción.
  Primero hay que ver que $\Gamma1 + \Gamma1$ es un subespacio de $A$. Como $\Gamma 1=1$ y tenemos elementos globales 
  globales $a$ y $b$, por la propiedad universal del coproducto, existe un morfismo $a+b\colon 1+1\to A$. Veamos que es un monomorfismo.
  \begin{equation*}
    \begin{tikzcd}
      (a+b)\times_A (a+b)\ar{r}\ar{d} & 1+1\ar{d}{a+b}\\
      1+1\ar{r}{a+b} & A.
    \end{tikzcd}
  \end{equation*}
  En un topos el coproducto es estable bajo productos fibrados, entonces $(a+b)\times_A (a+b)\cong
  (a\times_A a) + (a\times_A b) + (b\times_A a) + (b\times_A b)$. Además, $a\times_A b\cong 0$
  ya que $a\neq b$. Por lo tanto, $(a+b)\times_A (a+b)\cong (a\times_A a) + (b\times_A b)$, pero $a$ y $b$
  son monomorfismos, entonces $(a+b)\times_A (a+b)\cong 1+1$. Por lo anterior, $a+b$ es mono. Por otro lado,
  $\Gamma1+\Gamma1$ es un espacio discreto porque es el coproducto en $M$. Falta demostrar que los puntos
  de $A$ se factorizan a través de $\Gamma1+\Gamma1 = 1+1$.
  \begin{equation*}
    \begin{tikzcd}
      1\ar{r}{a}\ar[swap]{dr}{i} & \Gamma A & 1\ar{l}{b}\ar{dl}{i} \\
      {} & 1+1.\ar[dashed]{u}{a+b} 
    \end{tikzcd}
  \end{equation*}
  Los puntos de $A$ se encuentran en $\Gamma A$ y se factorizan a través de $a+b$. Por lo tanto,
  la unicidad de $\Gamma A$ implica que $\Gamma A = 1+1$.
\end{proof}

\begin{corollary}[label=cor:mbooleano]
  $M$ es un topos booleano.
\end{corollary}
\begin{proof}
  Sabemos por \ref{cor:bivaluado} que $\Omega$ tiene dos puntos, entonces por \ref{thm:puntitos} $\Gamma\Omega \cong 1+1$.
\end{proof}

\begin{theorem}[label=thm:modelocs]
    El topos $\toposM$ de espacios discretos es un modelo de CS.
\end{theorem}
\begin{proof}
    Ya demostramos que $M$ es un topos~\ref{thm:toposdiscretos}, satisface el Nullstellensatz pues es un axioma que le pedimos a $\esp$~\ref{ax:gds6} y
    es un topos booleano~\ref{cor:mbooleano}. Además, por~\ref{thm:onndecidible} $M$ hereda el ONN de $\esp$.
\end{proof}